\centerline{\bf \Huge Оценка + пример II}

\begin{flushright}
        \scriptsize{Никогда не думайте, что вы уже всё знаете. И как бы высоко ни оценивали вас, всегда имейте мужество сказать себе: «Я невежда».}
\end{flushright}

\problem{1}
Каким наибольшим количеством монет в 3 и 5 коп. можно набрать сумму в 1000 копеек?

\problem{2}
Какое наименьшее число клеточек на доске 8×8 можно закрасить в чёрный цвет так, чтобы была хотя бы одна закрашенная клетка в любом уголке из трёх клеточек?

\problem{3}
На какое наибольшее количество разных прямоугольников можно разрезать по линиям сетки прямоугольник 12×6 клеточек?


\problem{4}
На шахматной доске стоят 10 шахматных фигур – слоны и ладьи. Никакая фигура не бьёт ни одну другую. Какое наименьшее количество слонов может быть среди них? (ладья бьет по вертикали и горизонтали, слон по диагоналям)

\problem{5}
Какое наименьшее число выстрелов в игре "Морской бой" на доске $7\times7$ нужно сделать, чтобы наверняка ранить четырехпалубный корабль (четырехпалубный корабль состоит из четырех клеток, расположенных в один ряд)?


\problem{6}
В спортклубе тренируются 100 толстяков весом от 1 до 100 кг (у всех разный вес). На какое наименьшее число команд их можно разделить так, чтобы ни в одной команде не было двух толстяков, один из которых весит вдвое больше другого?

